\begin{lemma}[A source-mapped local slow graph]\label{wed:lem:slowgraph}
Property \textup{(F)} holds for system \eqref{wed:eq:fastslow}.
\end{lemma}
\begin{proof}
We apply the invariant-graph persistence theorem of Eldering
\cite[Theorem~3.2 and Remark~3.3(iii)]{Eldering13}, on a boundaryless
Euclidean base. The construction below checks its hypotheses and obtains
joint parameter regularity by adjoining a frozen, scaled parameter.

Fix $\mathsf K$ and $N$. Choose a bounded open $\mathsf O$ with
$\mathsf K\subset\mathsf O$ and a smooth compactly supported function
$\chi(U,P,r)$ equal to one on a neighbourhood of $\overline{\mathsf O}$.
Put $z=Q+g(U)$. In fast time $s=\tau/\varepsilon$, the transformed system is
\[
 U_s=\varepsilon P,\quad P_s=\varepsilon(z-g),\quad r_s=0,
 \qquad z_s=-z+\varepsilon gP-r\varepsilon^2g.
\]
Here $g'=g+1$ explains the cancellation in the last equation.
Choose a smooth compactly supported function $b(a)$ with $b(a)=a$ for
$|a|\le1$, and set $e(a)=\delta_*b(a)$, where $\delta_*>0$ will be small.
On $X\times Y=\R^4\times\R$, with base coordinates $(U,P,r,a)$ and
normal coordinate $z$, consider
\begin{align*}
 U_s&=\chi e(a)P,&P_s&=\chi e(a)(z-g),&r_s&=a_s=0,\\
 z_s&=-z+\chi\bigl(e(a)gP-r e(a)^2g\bigr).
\end{align*}
The unperturbed field is $(0,0,0,0,-z)$. Its zero section has identity
tangent flow and normal flow $e^{-s}$, with constant splitting and bounded
projections. It is therefore normally attracting with every finite
regularity gap. The base $X$ is connected, complete, and of bounded geometry.
The initial graph is zero, so the auxiliary graph family in Theorem~3.2
may be taken constant, by Remark~3.3(iii).

On any fixed tube $|z|<d$, both fields have bounded, uniformly continuous
derivatives of every prescribed finite order. Indeed, $\chi$ bounds all
occurrences of $e^U,P,r$, and $b$ bounds the frozen parameter. Their
difference tends to zero in every fixed finite $C^j$ norm as
$\delta_*\to0$. Apply Theorem~3.2 with $j=\max(N,2)$ and H\"older index
zero. For sufficiently small $\delta_*$ it supplies a $C^j$ invariant
graph $z=h(U,P,r,a)$ over $X$. This application uses neither the book's
boundary extension nor an overflowing-boundary hypothesis.

At $a=0$, the base velocity is zero and $z_s=-z$. Tangency forces
$h(U,P,r,0)=0$. For $0\le\varepsilon\le\varepsilon_0:=\delta_*/2$ put
\[
 w(U,P,r,\varepsilon)
   =h(U,P,r,\varepsilon/\delta_*)-g(U).
\]
This function is jointly $C^N$ on a neighbourhood of
$\overline{\mathsf O}\times[0,\varepsilon_0]$, and $w(U,P,r,0)=-g(U)$.
Over $\mathsf O$, $\chi=1$ and $e(\varepsilon/\delta_*)=\varepsilon$,
so graph tangency gives exactly
\[
 \varepsilon(w_UP+w_Pw)=-w-\varepsilon P-(1+r\varepsilon^2)g(U).
\]
This is \eqref{wed:eq:invariance}, including at $\varepsilon=0$.
\end{proof}

The graph and the smallness threshold may depend on the compact set and on
$N$. The lemma does not give one $C^\infty$ graph on a fixed parameter
interval, an explicit threshold, or uniform control on growing energy
balls. It establishes local invariance for the original system, not global
invariance of the cutoff graph under that system. Theorem~3.2, its relevant
remarks and the author's errata were inspected for this application; the
historical Fenichel/Jones numbering is not needed.
